% ============================================================
%  Chapter: Values, Ethics and Human Interface
%  Foundation Batch --- Lectures 1 \& 2
%  Topics: Values \textbullet\ Ethics vs Morals \textbullet\ Ethics vs Law
%          \textbullet\ Constitutional Morality \textbullet\ Determinants of Ethics
% ============================================================

\chapter{Ethics and Human Interface}
\vspace{14pt}

\section{Values}

\subsection{Meaning of Values}

When we say something is ``valuable'', we mean that it is important to us. To call a pen or a phone valuable is to attach a degree of importance or significance to it --- perhaps because it is costly, or because there is emotion attached to it.

\begin{KeyConcept}
\begin{itemize}
\item Values are the degree of importance and significance that we attach to something.
\item As human beings we prefer different things, and each preference reflects a value.
\end{itemize}
\end{KeyConcept}

\begin{ExampleBox}
\begin{itemize}
\item Raghukul (Ramayana): ``Raghukul reet sada chali aayi, praan jaaye par vachan na jaaye'' --- one must uphold one's promise even at the cost of one's life. Upholding promises is therefore a core value of Raghukul, because that is what they find most important.
\item Tribals find the forest valuable --- not merely for economic livelihood, but as a way of life. They are culturally dependent on it, often believe their ancestors live in it, and derive immense solidarity from it. The forest is central to their entire existence.
\item For Gandhiji, non-violence was supremely important. When the Chauri Chaura incident occurred, he withdrew the Non-Cooperation Movement because independence itself was not desirable to him at the cost of violence --- non-violence was clearly the more preferable value.
\end{itemize}
\end{ExampleBox}

\subsection{Types of Values --- Moral, Immoral and Amoral}

A preference can carry a moral dimension, an immoral dimension, or no moral dimension at all. This gives three categories of values.

\paragraph{Moral values} Moral values are those preferences that, by widely held standards, are considered right --- for example honesty and loyalty. Satyavadi Harishchandra, who preferred truth above all, embodies a moral value.

\paragraph{Immoral values} Immoral values are preferences generally regarded as wrong. A greedy person prefers greed above everything, so greed becomes that person's value --- an immoral one. For a rapist, lust becomes more important than human dignity, which is again an immoral value.

\begin{CriticalPoint}
When we indulge in inhuman acts --- the Madhya Pradesh incident of urinating on another person, imputing a person's limb as ``justice'', bonded labour, manual scavenging, rape, or human trafficking --- we fail to treat other human beings as human beings. In such actions we are bereft of human values.
\end{CriticalPoint}

\paragraph{Amoral values} Amoral values are preferences that carry no connotation of righteousness or wrongness. Liking a particular singer, colour, clothing or music is an amoral preference --- you cannot be called wrong for preferring blue, or for liking Arijit Singh.

\begin{ExampleBox}
Preferring the colour white over black is an amoral choice. But saying ``I love white people and hate black people'' is immoral, because it introduces race and denies the equal humanity of others.
\end{ExampleBox}

\begin{KeyConcept}
Attitude, like values, can also be moral, immoral or amoral. An attitude that carries no connotation of right or wrong --- such as one's attitude towards a pen --- is an amoral attitude.
\end{KeyConcept}

\subsection{Relative Hierarchy of Values}

A human being holds many values at once. You may value time spent with loved ones, your career ambition of becoming a civil servant, and your own health --- all simultaneously, and all as part of a single value system.

\begin{ExampleBox}
Prelims is tomorrow, and your best friend arrives insisting on a party. You value your friend, and you also value becoming an IAS officer. Both are genuine values, and here they collide.
\end{ExampleBox}

\begin{KeyConcept}
\begin{itemize}
\item Values have an internal, relative hierarchy --- we prefer some over others.
\item This hierarchy is contextual and dynamic; it keeps changing with circumstances.
\item At present, becoming a civil servant may be the top priority, so that you are temporarily unable to give full time to health or to loved ones.
\end{itemize}
\end{KeyConcept}

\subsection{Ethical Dilemmas Arising from Value Conflict}

An ethical dilemma occurs when you are unable to decide between two competing ethical claims --- where both options are ethical, yet you cannot pursue both.

\begin{ExampleBox}
\begin{itemize}
\item Sunil, a young and dynamic SP, is posted in a district where illegal sand mining is rampant. When he acts against it, the mining mafia threatens to harm his family. On one side is his responsibility towards his family; on the other is his duty and integrity. Both are ethical values, and the competing claims create a dilemma. Such dilemmas are resolved through ethical reasoning (utilitarianism and other frameworks, to be studied later).
\item COVID warriors continued their duty knowing that anyone's life could be at risk. Their choice reflects selflessness and the public-service philosophy --- upholding one's duty even at great personal cost.
\item Bhishma Pitamah (Mahabharata) took an oath to serve Hastinapur faithfully and held to it because the oath was tied to his very identity; for him it was the most important thing.
\end{itemize}
\end{ExampleBox}

\begin{QuickRevision}
\begin{itemize}
\item Individuals hold multiple obligations and value systems at once.
\item When their relative hierarchy is clear, we act; when it is not, a dilemma arises.
\item Dilemmas are resolved by ethical reasoning, developed through case studies.
\end{itemize}
\end{QuickRevision}

\section{Collective Values}

Values are not only individualistic; they are also collective. Great individuals --- Daanvir Karna, Bhishma Pitamah, Satyavadi Harishchandra, and Gandhiji (non-violence, Satyagraha) --- are known by the values they upheld throughout their lives, and those values guided their decision-making. But preferences held in common by a group are collective values.

\subsection{Family Values}

Families play a huge role in the inculcation of values, so many of a person's standards are shaped from childhood. Discipline can be a collective value of police or army families; respecting elders is a collective family value; and patriarchal families collectively hold patriarchal values.

\subsection{Societal Values}

Different societies hold different collective values, and these shape everyday behaviour.

\begin{center}
\begin{tabularx}{\textwidth}{>{\raggedright\arraybackslash}p{3.4cm} Y Y}
\rowcolor{AccentDark}
\thead{Attribute} & \thead{Urban society} & \thead{Rural society} \\
\toprule
Dominant value & Individualism & Collectivism \\
Sense of belonging & Weaker; people think of themselves & Stronger; ``our own people'' \\
Response to a road accident & Fewer people stop to help & More likely to help --- higher social capital \\
\bottomrule
\end{tabularx}
\end{center}

\begin{ExampleBox}
\begin{itemize}
\item Japanese society: discipline is a collective value --- civic sense lets huge numbers cross intersections without a jam; cleanliness is another collective value.
\item Okinawa community (Japan): many people average more than 100 years of age, and one stated reason is the strong sense of belongingness in the community, which is treated as an extension of one's own life.
\item Scandinavian / Nordic countries: trust is the collective value system --- items left in a train or washroom remain safe. Trust-based societies see crime fall; privacy is also treated as very important there.
\end{itemize}
\end{ExampleBox}

\begin{CriticalPoint}
Suspicion-based societies rely on instruments like speed cameras on highways; trust-based societies need fewer such controls, and crime declines as a result.
\end{CriticalPoint}

Many social problems are driven by collective societal values rather than by government alone --- dowry, corruption, air pollution, paper leaks, and the criminalisation of politics. Corruption, for instance, has a demand side rooted in society, not only a supply side in government.

\begin{PYQLink}
``Corruption is the manifestation of the failure of core values of society.'' A major reason corruption persists is the failure of societal values --- we glorify those who have prospered by wrong means, engage in hero-worship of criminals, and celebrate the wrong role models in films and popular culture.
\end{PYQLink}

\begin{ExampleBox}
Negative societal values in India: roughly 23\% of girls are married before 18 years of age, showing that child marriage remains prevalent even though it is illegal; domestic violence, though ethically and legally wrong, is still normalised in places.
\end{ExampleBox}

\subsection{Constitutional Values}

The Constitution too has a collective value system. It collectively prefers and attaches importance to justice, liberty, equality and fraternity, and to being sovereign, socialist, secular, democratic and republic.

\begin{KeyConcept}
\begin{itemize}
\item These constitutional values are moral values, collectively preferred by the Constitution.
\item Taken together they constitute constitutional morality --- the spirit of the Constitution, which lies much beyond its mere letters.
\end{itemize}
\end{KeyConcept}

\subsection{Organizational \& Professional Values}

\begin{center}
\begin{tabularx}{\textwidth}{>{\raggedright\arraybackslash}p{3.0cm} Y Y}
\rowcolor{AccentDark}
\thead{Sphere} & \thead{Collective value system} & \thead{Consequence for work culture} \\
\toprule
Private sector & Profit-making, efficiency, confidentiality & Faster decision-making \\
Public sector & Public interest & Slower decisions --- accountability, public money, elaborate checks and balances \\
Medical profession & Hippocratic oath --- saving the patient's life & Duty to preserve life \\
ISRO & Frugality, innovation & Low-cost missions (Chandrayaan / Mangalyaan cheaper than a film budget) \\
\bottomrule
\end{tabularx}
\end{center}

Values thus operate at the level of the individual, the family, society, the Constitution, organizations, and professions. ISRO's frugality is even treated as a national value, just as an entrepreneur may stand for the value of entrepreneurship.

\begin{ExampleBox}
\begin{itemize}
\item Tribal societies hold collective values of egalitarianism, solidarity and contentment (santosh). ``Tribe'' means a group of people with a common characteristic; the belief that all members descend from a common ancestor produces a strong sense of solidarity.
\item Communities may hold wildlife management as a collective value.
\item Societies may collectively value materialism (consumerism) or spiritualism.
\end{itemize}
\end{ExampleBox}

\begin{QuickRevision}
\begin{itemize}
\item Values operate at individual, family, societal, constitutional, organizational and professional levels.
\item Urban societies lean individualistic; rural and tribal societies lean collective (solidarity, egalitarianism).
\item Collective societal values explain many public problems (corruption, dowry, pollution) and their solutions.
\item Constitutional morality is the collective moral value system of the Constitution --- its spirit beyond its letter.
\end{itemize}
\end{QuickRevision}

\section{Meaning of Ethics}

Having understood values --- the degree of preference and importance attached to something --- we turn to ethics itself. When we judge a human behaviour as right or wrong, we do so against some standard. If we fail to perform our duty and call that behaviour wrong, then duty is functioning as the ethical standard.

\begin{KeyConcept}
\begin{itemize}
\item Ethics = the standards on the basis of which we evaluate human actions as right or wrong.
\item For example, media ethics are the standards applied to the media --- such as the requirement that information be verified before it is published.
\end{itemize}
\end{KeyConcept}

\subsection{Sources of Ethics}

Ethical standards --- the source of our judgment of right and wrong --- are shaped by several factors.

\begin{itemize}
\item Religion: each religion prescribes principles. Jainism emphasises non-violence; Buddhism teaches the middle path and that desire is the root cause of sorrow; Hindu, Muslim and Christian traditions each give their own ethical principles.
\item The Gita --- Nishkaam Karma: act without attachment to the consequences or rewards; perform your duty as a duty, because it is the right thing to do.
\item Culture: the culture we are born into tells us what is right and wrong --- for example, the norms of the Bishnoi community, whose members would regard a violation of their conservation ethic as unethical.
\item Law: law also prescribes standards of right and wrong.
\end{itemize}

Thus the ethical judgment of right and wrong is shaped by many factors together --- culture, religion and law among them.

\section{Ethics vs Morals}

Both ethics and morals are standards that tell us what is right and what is wrong. The distinction lies in their source.

\begin{KeyConcept}
\begin{itemize}
\item Ethics = standards that are external to us --- those prescribed to us by law, family, custom, culture or religion.
\item Morals = standards that are intrinsic to us --- our own personal judgment of what is right and wrong.
\end{itemize}
\end{KeyConcept}

\begin{ExampleBox}
\begin{itemize}
\item If your family tells you ``do not steal'', that is an ethical standard (externally prescribed). If you yourself also believe stealing is wrong, that is your moral standard. Here ethics and morals coincide.
\item If your family tells you not to marry inter-caste but you believe such marriage is perfectly right, then your morals diverge from the externally prescribed ethical standard, and the two come into conflict.
\end{itemize}
\end{ExampleBox}

\begin{center}
\begin{tabularx}{\textwidth}{>{\raggedright\arraybackslash}p{2.6cm} Y Y}
\rowcolor{AccentDark}
\thead{Basis} & \thead{Ethics} & \thead{Morals} \\
\toprule
Source & External --- prescribed to us & Internal --- our own judgment \\
Subjectivity & More uniform / objective & More subjective; varies person to person \\
Enforcement & Usually backed by external sanctions (religious, social, professional, organizational) & Rely on internal guilt, shame and regret \\
\bottomrule
\end{tabularx}
\end{center}

\begin{CriticalPoint}
Because ethical standards are externally applied, they tend to be relatively uniform for a group and come with sanctions --- a professional body or organization can punish a breach. Moral standards, being intrinsic, are enforced only by one's own conscience.

A whistleblower who exposes wrongdoing despite the risk of being fired is typically driven by moral sanction --- the inability to carry the guilt of staying silent while people are harmed.
\end{CriticalPoint}

\subsection{When Ethics and Morals Conflict}

Ethics and morals can go hand in hand when you believe in the standards prescribed to you; they conflict when your personal morality diverges from those external standards.

\begin{ExampleBox}
\begin{itemize}
\item A lawyer defends a client and secures a fair trial because professional ethics demand representation, even if the lawyer personally believes the client is guilty.
\item A doctor is bound by medical ethics to save a patient's life --- even a criminal's --- regardless of personal feeling.
\item A vegetarian appointed chairman of a fishery board faces conflict when professional ethics demand promoting fishing against a personal moral stance.
\item A Christian may face conflict when religion (external ethics) holds abortion to be wrong, while her personal morality asserts bodily autonomy --- for instance where the pregnancy resulted from rape.
\item Untouchability was considered ethical by society under the purity--pollution principle, but reformers such as Dr B.R. Ambedkar knew intrinsically it could never be right and argued against it.
\item A student's own morality may be against wearing the hijab while society or the government prescribes it --- ethics and morals again collide.
\end{itemize}
\end{ExampleBox}

\begin{GovtPolicy}
Guru Nanak emphasised ``kirat'' --- doing one's work honestly --- and upheld the dignity of labour against the idea that certain occupations are polluting or menial. The notion of untouchability sprang precisely from the purity--pollution belief that some jobs were impure.
\end{GovtPolicy}

\section{Ethics vs Law}

\begin{PYQLink}
The interface of ethics and law is one of the most important themes of the paper --- a question on this theme appears almost every year.
\end{PYQLink}

Both ethics and law tell us what is the right thing to do. Law tells us not to indulge in corruption, not to take dowry, that domestic violence and triple talaq are wrong. Ethics tells us the same. The differences lie elsewhere.

\begin{center}
\begin{tabularx}{\textwidth}{>{\raggedright\arraybackslash}p{2.6cm} Y Y}
\rowcolor{AccentDark}
\thead{Basis} & \thead{Law} & \thead{Ethics} \\
\toprule
How it is made & Follows a due procedure / due process & No corresponding due process; not every ethical standard becomes law \\
Subjectivity & More objective and uniform --- all governed by the same law & More subjective --- different societies hold different standards \\
Enforcement & Formal institutions --- police, judiciary, prison & Lacks formal institutions of enforcement \\
Breadth & Narrower --- minimum morality & Broader --- maximum morality \\
\bottomrule
\end{tabularx}
\end{center}

\begin{KeyConcept}
Law is minimum morality; ethics is maximum morality. Not every ethical standard is legally enforced.
\end{KeyConcept}

\begin{ExampleBox}
\begin{itemize}
\item Not all fundamental duties are legally enforced: not voting is generally regarded as unethical, yet a citizen cannot be punished for it.
\item Marital rape is not illegal --- one cannot be tried for it in court --- yet it is unethical, because consent is absent.
\end{itemize}
\end{ExampleBox}

\subsection{Legal but Unethical}

Because law is only minimum morality, many actions are perfectly legal yet remain unethical, since the actor is unaware of, or ignores, a higher set of ethical obligations.

\begin{itemize}
\item Emergency in India: declared legally and constitutionally, with all procedures followed, yet unethical because it violated the democratic spirit.
\item Abuse of discretion: a Governor who sits indefinitely on a bill, a Speaker who misuses anti-defection powers to manufacture a government, or the passing of a bill as a money bill to bypass Rajya Sabha scrutiny --- each is technically legal but grossly unethical and anti-democratic, because an elected government's mandate must be respected.
\item Bulldozer justice: demolishing a house citing violation of construction norms is technically legal, but selective demolition of only one house while similar violations nearby are ignored is grossly unethical. The Supreme Court has questioned exactly this selective application --- ``show me the man, I will show you the rule.''
\item Surrogate advertisements: brands such as Kingfisher and Imperial Blue advertise ``drinking water'' or soda to bypass the ban on alcohol advertising --- legal, but unethical.
\item Road-accident compensation: tribunals compute compensation by factoring in the income the deceased was earning, effectively putting a price tag on human life. This is legal but unethical --- it treats human beings merely as a means, not as an end in themselves, and implies one life is worth more than another.
\item Parliamentary practice: fewer sitting days and referring bills to committees only rarely is legal but unethical, since it weakens parliamentary accountability.
\end{itemize}

\begin{DataFact}
\begin{itemize}
\item In India, less than a fifth --- around 20\% --- of bills are referred to a parliamentary committee.
\item By contrast, in the UK bills are examined by committees of both Houses.
\end{itemize}
\end{DataFact}

\begin{itemize}
\item Bureaucratic red-tapism: using rules as an instrument to harass individuals, with rules becoming ends in themselves, is legal but grossly unethical --- it violates the dignity of the citizen.
\end{itemize}

\begin{ExampleBox}
\begin{itemize}
\item T.N. Seshan used the very same rules and powers of the Election Commission --- powers that were in fact reduced when it became a multi-member commission --- in a highly positive manner to enhance governance. The Supreme Court observed that the ECI's history can be divided into a pre-Seshan and a post-Seshan phase. The lesson: a positive-minded officer uses the same rules to improve governance, while a negative-minded officer misuses them to harass citizens.
\item PM's moral appeals: violating a moral appeal (for example on giving up a subsidy) may be perfectly legal, yet unethical in the contemporary context.
\end{itemize}
\end{ExampleBox}

\subsection{Illegal but Ethical}

The converse also exists: some actions are illegal yet ethical, often where the law itself is unethical.

\begin{itemize}
\item Civil disobedience: Gandhiji broke the salt law --- illegal, but ethical, because the law itself was not ethical.
\item Abortion: abortion is banned in several US states, so performing it there may be illegal, yet it can be ethical.
\end{itemize}

\begin{GovtPolicy}
Unbanked Direct Blood Transfusion (UDBT): in India, civilians must source blood legally from recognised blood banks, and direct person-to-person transfusion is not permitted (though it is allowed for army personnel). Yet in remote areas and emergencies --- road accidents, or heavy blood loss during childbirth that can be fatal --- blood is screened for basic compatibility (blood group, STDs, hepatitis B) and transfused directly to save a life. This is illegal but ethical. The Drug Technology Advisory Board (DTAB) was asked to allow it, but there is a countervailing concern about controlling the transmission of disease.
\end{GovtPolicy}

\subsection{Right vs Right}

\begin{PYQLink}
``Ethics is knowing the difference between what you have a right to do and what is the right thing to do.'' (PYQ.)
\end{PYQLink}

Here the ``right to do'' is your entitlement, while the ``right thing to do'' is the ethically correct action. The two can diverge.

\begin{ExampleBox}
\begin{itemize}
\item Students have every right to keep asking doubts, but constantly doing so is not the right thing, as it disrupts the pace of the class and everyone's learning.
\item You have the right not to help a road-accident victim --- no one can force you --- but helping is the right thing to do.
\item A state government has the right to recover a farm loan, but during a drought that may not be the right thing to do, which is why farm-loan waivers are granted.
\item The right to seek information under RTI can be misused by filing excessive requests to harass an office --- an entitlement exercised in a way that is not the right thing to do.
\end{itemize}
\end{ExampleBox}

\section{Constitutional Morality}

When one is caught between ethics and law, the guiding concept is constitutional morality.

\begin{KeyConcept}
Constitutional morality refers to the principles and values enshrined in the Constitution --- not just the letter but the spirit of the Constitution; not just its literal interpretation but its underlying spirit.
\end{KeyConcept}

\begin{ExampleBox}
\begin{itemize}
\item Maneka Gandhi case: earlier the test was merely ``procedure established by law'' --- whether an action complied with the procedure. But the procedure itself may not be right, so the Court moved to the principles of natural justice and ``due process of law'' --- the spirit of the Constitution.
\item Article 21: its literal meaning is very narrow, yet the courts have expanded it enormously --- to environment, housing and the right to privacy --- by going by the spirit of the Constitution. The Sabarimala judgment is another instance of reasoning from spirit rather than letter.
\end{itemize}
\end{ExampleBox}

\begin{GovtPolicy}
\begin{itemize}
\item RTI Act: the provision that Public Information Officers must render ``reasonable assistance'' to applicants can be read narrowly (``the website is there, check it'') or, by an officer guided by constitutional morality, broadly --- breaking linguistic and financial barriers and actively helping the applicant file the request.
\item Article 19: imposing Section 144 merely to silence criticism would be following the letter, not the spirit, of ``reasonable restriction''.
\end{itemize}
\end{GovtPolicy}

\begin{CriticalPoint}
The letter of the Constitution is constitutional law; constitutional morality is the wider ethics --- the spirit. Decision-making should be guided by the spirit, not merely the letter.
\end{CriticalPoint}

Constitutional morality is a collective responsibility of all stakeholders, not of government alone. Citizens uphold it by discharging their fundamental duties and the ideal of fraternity --- avoiding hate speech and caste-based conflict; parents uphold it through their duties (Article 51A); and the State upholds it through its own conduct.

\begin{QuickRevision}
\begin{itemize}
\item Constitutional morality = following the spirit, not merely the letter, of the Constitution.
\item It is upheld collectively by citizens, families and the State.
\item In any law--ethics interface answer, conclude that constitutional morality should guide decision-making --- having the right to do something does not mean one should do it.
\end{itemize}
\end{QuickRevision}

\section{Determinants of Ethics}

\begin{KeyConcept}
\begin{itemize}
\item Determinants of ethics are the factors that help in both the formation and the evolution of ethical standards.
\item Ethical standards do not remain static; they also evolve and change over time.
\end{itemize}
\end{KeyConcept}

\subsection{Institutional \& Human Determinants}

\begin{itemize}
\item Religion: different religions emerged across human history, each adding ethical standards --- Buddhism (the middle path; desire as the root cause of sorrow) and Jainism (Anekantavada --- multiple versions of truth). This is why questions ask how the teachings of Buddhism, Mahavira or Jainism remain relevant today: their ethical standards continue to guide humanity.
\item Legislature: by passing laws it changes the meaning of right and wrong --- online betting apps that were legal have been made illegal; triple talaq, once accepted, has been rendered unjust.
\end{itemize}

\begin{GovtPolicy}
India became the first country to legally mandate CSR through the Companies Act, 2013 --- a clear example of the legislature acting as a determinant of ethics.
\end{GovtPolicy}

\begin{itemize}
\item Executive and Judiciary: through orders and judgments they reshape our perspective on right and wrong. The Maneka Gandhi case shifted the norm from ``procedure established by law'' to ``due process of law''; the Kesavananda Bharati case established that all parts of the Constitution can be amended except the basic structure.
\item Individuals (self-reflection): Ashoka, after the Kalinga war, turned to peace and propagated it through the Ashokan pillars, changing society's ethical outlook; Raja Ram Mohan Roy campaigned against Sati Pratha and thereby changed the ethical standards of society.
\item Society --- family and culture: the family and culture we are born into shape our standards. Cultural norms (as with the Bishnoi community) are hard to change because we are status-quoist by nature.
\item Professional institutions: media ethics, engineering ethics and medical ethics each shape conduct within a profession.
\end{itemize}

\subsection{Technology as a Determinant}

Technology plays a huge role in reshaping ethical standards.

\begin{ExampleBox}
\begin{itemize}
\item AI and copyright --- the Anthropic case: a US court ruled that training AI on books without permission amounted to fair use, because the model creates something substantially new (a new common good) rather than reproducing the original. However, the storage of pirated books in a library constituted copyright infringement, and the company was held liable for that. Thus the meaning of copyright is itself shifting because of technological innovation.
\item The EU's GDPR regime protects personal data as intrusive technology grows.
\item The UK is moving to introduce laws against AI tools used to generate sexual-abuse images --- new technology forcing new ethical and legal standards.
\end{itemize}
\end{ExampleBox}

\subsection{Economic \& Global Determinants}

\begin{itemize}
\item Economic models: the rise of the gig economy and informal work is changing older norms. Gujarat amended its Factories Act to allow women to work night shifts and longer hours --- what was once considered just is reinterpreted as new economic models emerge.
\item Globalization and multilateral frameworks: India moved from process patent to product patent after signing the TRIPS agreement. A process patent protects only the method of making a product (weaker protection); a product patent protects the product itself (stronger protection). Changing global norms on intellectual property made what was earlier just now unjust.
\item Nagoya Protocol --- Access and Benefit Sharing: benefits must be shared with the community whose resources are used (detailed treatment belongs to environment, not ethics).
\item Thinkers: during the French Revolution, when atrocities were justified in the name of religion, Voltaire said that God lives in the human heart and not in churches --- changing people's thinking and ethical standards. Revolutions use literature, media and the press to spread new ideas such as humanism.
\end{itemize}

\subsection{Evolution of Norms --- The Concept of Justice}

The clearest illustration of evolving ethical standards is the concept of justice, which has moved from retribution to rehabilitation.

\begin{center}
\begin{tabularx}{\textwidth}{>{\raggedright\arraybackslash}p{2.9cm} Y Y}
\rowcolor{AccentDark}
\thead{Aspect} & \thead{Earlier understanding} & \thead{Present understanding} \\
\toprule
Aim of justice & Retribution --- punishment-oriented & Rehabilitation and reformation \\
Treatment of offenders & Imputing a limb was seen as justice & Rejected --- even dreaded criminals have human rights, a fair trial and dignity \\
Guiding code & IPC --- Indian Penal Code (penalty) & BNS --- Bharatiya Nyaya Sanhita (nyaya / complete justice, attitude change) \\
\bottomrule
\end{tabularx}
\end{center}

\begin{CriticalPoint}
The same action --- mutilating a criminal --- was considered just in the past but is unjust today, because our understanding of justice has evolved to include human rights. Even a dreaded leader's body is buried as per norms; everyone deserves a fair trial and protection of dignity. This shift is reflected in community service as a form of punishment under the new code.
\end{CriticalPoint}

Similar evolution is visible elsewhere: the expansion of human rights and the recognition of the right to privacy under Article 21 (driven partly by intrusive technology, hence GDPR); and the recognition of LGBT rights --- once considered unethical, now just --- as our understanding of gender and sexual orientation has evolved.

\subsection{Changing Context Must Be Under Constant Scrutiny}

\begin{PYQLink}
``What was just some time back may turn out to be unjust.'' Changing context should therefore be kept under constant scrutiny.
\end{PYQLink}

\begin{ExampleBox}
\begin{itemize}
\item Copyright standards must be re-examined as technology changes, or they lag behind reality.
\item Early maturity --- Juvenile Justice Act: earlier no person below 18 could be tried as an adult even for a heinous crime; recognising that maturity may be attained before the biological age of 18, the law now allows a juvenile to be tried as an adult in heinous cases, so that such offenders do not go scot-free --- otherwise it would be a travesty of justice.
\item Romeo--Juliet clause --- POCSO Act: since there is no concept of consensual sex where one partner is a minor, genuine relationships can be caught by the law, and the provision is misused by families to target the boy. Courts have observed that changing context must be understood to prevent a miscarriage of justice.
\item Facial tattooing: tribal girls once received facial tattoos to appear unattractive and avoid sexual exploitation by landlords and others during and before British rule; with that context gone, the standard around the custom changes.
\item Officials performed the last rites of Baswaraju --- reflecting the evolution of the value of human dignity.
\end{itemize}
\end{ExampleBox}

\begin{QuickRevision}
\begin{itemize}
\item Determinants are all the factors that form and evolve ethical standards.
\item Institutions (religion, legislature, executive, judiciary, professions), individuals, society and culture all contribute.
\item Technology, economic models, globalization and thinkers are key contemporary determinants.
\item Norms evolve --- justice from retribution to rehabilitation, expansion of human rights and privacy --- so changing context must be scrutinised constantly.
\end{itemize}
\end{QuickRevision}
